\section{The two-moment stabilizer theorem}

Let $K_p=\Q(\zeta_p)$.  For $a\in\F_p^\times$, write
$\sigma_a(\zeta_p)=\zeta_p^a$.

\begin{lemma}[Galois--histogram dictionary]\label{lem:dictionary}
One has
\[
 \sigma_a(B_{p,1})=B_{p,1}
 \quad\Longleftrightarrow\quad
 H_p(ar)=H_p(r)\quad\text{for all }r\in\F_p.
\]
\end{lemma}

\begin{proof}
After relabeling, the coefficient vector of
$\sigma_a(B_{p,1})-B_{p,1}$ is
$H_p(a^{-1}r)-H_p(r)$.  A rational polynomial of degree at most $p-1$
that vanishes at $\zeta_p$ is a scalar multiple of
$\Phi_p(X)=1+X+\cdots+X^{p-1}$.  The coefficient sum is zero because both
histograms have total mass $2p^2$.  The scalar is therefore zero.  The
converse is immediate.
\end{proof}

Write $p-1=3m$.  Since $p\equiv1\pmod6$, $m$ is even.  For even $k$, define
\[
 M_k=\sum_{r\in\F_p}r^kH_p(r)
     =2\sum_{x,y\in\F_p}f_p(x,y)^k\quad\text{in }\F_p.
\]

\begin{proposition}[Two exact power moments]\label{prop:moments}
For every split prime,
\[
 M_{2m}=2\binom{2m}{m}4^m\ne0\pmod p.
\]
For $p\ge13$,
\[
 M_{2m+2}
 =2\frac{(2m+2)!}{(m-2)!^2\,6!}
   2^{2m-4}c^6\ne0\pmod p.
\]
\end{proposition}

The proof is given in \cref{sec:proofs}.  Its key feature is uniqueness: at
each exponent, precisely one multinomial term survives finite-field monomial
orthogonality.

\begin{theorem}[Histogram stabilizer]\label{thm:stabilizer}
For every prime $p\equiv1\pmod3$,
\[
 \Stab(H_p)=\{a\in\F_p^\times:H_p(ar)=H_p(r)\ \forall r\}
             =\{\pm1\}.
\]
\end{theorem}

\begin{proof}
If $a$ stabilizes $H_p$, then a change of variables in $M_k$ gives
$(a^k-1)M_k=0$.  For $p\ge13$, \cref{prop:moments} implies
$a^{2m}=a^{2m+2}=1$, hence $a^2=1$.  At $p=7$, the first formula gives
$M_4=3\ne0$, while every $a\in\F_7^\times$ satisfies $a^6=1$; again
$a^2=1$.  Conversely, $H_p$ is even, so both signs stabilize it.
\end{proof}
